% !TEX root = main.tex
\section{Generic Construction of \SMS}\label{section:generic-sms}

With the \SMS notion established, we build a generic scheme from standard primitives,
treating the hidden-circuit NIZKAoK as an abstract component
instantiated in Section~\ref{section:proving-function-evaluations-via-plonk}.
Following Nguyen~\etal~\cite{CRYPTO:NGSY22},
each enrolled user holds an $\FA$-issued certificate
$\certattr = (\attr, F, \vkattr, \sigma)$
binding their attribute, hidden policy, and a per-user verification key,
where $\sigma \gets \SSign(\mfk, (\attr, F, \vkattr))$
and the global verification key $\vk$ is embedded in $\gpk$.

At signing, the user encrypts the opening tag $\op$ under $\TA$'s public key $\ek$,
producing a ciphertext $\ctop$ that $\TA$ can later decrypt to recover $\op$.
A NIZKAoK proof $\PiGenericSMS$ certifies correct evaluation of $P(F(\cdot), G(\cdot))$,
validity of the membership certificate,
and correct encryption of $\op$,
without revealing $\attr$, $F$, or any intermediate computation trace.
How the NIZKAoK handles the hidden evaluation of $F$
via PLONK's arithmetization is developed in Section~\ref{section:proving-function-evaluations-via-plonk};
here, we treat the NIZKAoK as an abstract component for arithmetic circuit satisfiability.
%
In the \SMS instance for the CBDC example,
$F$ checks the spending limit, $G$ validates the merchant,
and $P$ sets $\op$ to the pseudonym or identity tag.

Without further binding, an adversary could modify $(\ctop, \PiGenericSMS)$
to produce a new valid \SMS signature.
%Following~\cite{PKC:BelFuc14,CRYPTO:NGSY22,ACISP:TNLPS24},
The user generates a one-time key pair $(\ovk, \osk)$ at the time of signing,
certifies $\ovk$ under $\skattr$,
and computes a one-time signature $\sigmaonetime$ over all signature components.
The NIZKAoK then proves validity of $\sigmaonetime$ under $\vkattr$
rather than proving the secret key $\skattr$ directly,
binding all components together against post hoc modification.

% !TEX root = main.tex
\subsection{Construction}\label{subsection:generic-SMS-description}
Suppose $\SignatureScheme = (\SKeyGen, \SSign, \SVerify)$ is
a secure signature scheme,
$\OTSScheme = (\OTSKeyGen, \OTSSign, \OTSVerify)$ is
a strongly unforgeable one-time signature scheme,
$\EncryptionScheme = (\EKeyGen, \EEnc, \EDec)$ is 
a CPA-secure PKE scheme with randomness space $\RandomnessEncrypt$.
We also use a NIZKAoK system $\NIZKAoK$
%$ = (\Setup, \Prove, \Verify)$ 
for arithmetic circuit satisfiability.
The message space $\MessageSpace$, derived message space $\TMessageSpace = \FF^\ellmsg$,
attribute space $\AttrSpace$, witness space $\WitnessSpace$,
function families $\FFamily$ and $\GFamily$, combiner $P$,
opening set~$\OpeningSet$, and signable subset $\OpeningSetSignable \subsetneq \OpeningSet$
are determined as part of $\SMSSetup$.

All computations are over $\FF$ (instantiated as $\ZZq$ for a prime $q$).
Each $F \in \FFamily$ and $G \in \GFamily$ share output length $\nout$,
so the combiner has type $P \colon \FF^{2\nout} \to \TMessageSpace \times \OpeningSet$.
The witness decomposes as $\bw = (\bw_F \| \bw_G)$,
where $\bw_F$ and $\bw_G$ support the computations of $F$ and $G$, respectively.

\begin{description}

\item[$\SMSSetup(1^\lambda, N) \to (\pp, \State, \mfk, \tsk)$:]
Sample $(\ek, \dk) \gets \EKeyGen(1^\lambda)$ for $\TA$,
$(\vk, \sk) \gets \SKeyGen(1^\lambda)$ for $\FA$,
and lastly 
$\ppnizk$ as the public parameter of $\NIZKAoK$.
Set $\gpk \coloneqq (\ek, \vk, \ppnizk)$, $\mfk \coloneqq \sk$, $\tsk \coloneqq \dk$, and $\State \coloneqq \emptyset$.
Determine $\MessageSpace$, $\TMessageSpace$, $\AttrSpace$, $\WitnessSpace$,
$\FFamily$, $\GFamily$, $P$, $\OpeningSet$, $\OpeningSetSignable$.
Return
\[
\pp \coloneqq (\gpk, \MessageSpace, \TMessageSpace, \AttrSpace,
\WitnessSpace, \FFamily, \GFamily, P, \OpeningSet, \OpeningSetSignable),\
\State,\
\mfk,\
\tsk.
\]

\item[$\SMSJoin\interaction{\User(1^\lambda)}{\FA(\mfk)}(\pp, \State, \attr, F)
\!\to\! (\State', (\certattr, \secattr))$:]
For user enrollment,
user $\User$ 
and 
function-issuing authority $\FA$
proceed as follows.

\begin{enumerate}
\item 
$\User$ computes $\sigmaenroll \gets \SSign(\skattr, (\attr, F, \vkattr))$,
where $(\vkattr, \skattr) \gets \SKeyGen(1^\lambda)$,
and sends $(\vkattr, \sigmaenroll)$ to $\FA$.
\item
$\FA$ aborts if $\SVerify(\vkattr, (\attr, F, \vkattr), \sigmaenroll) = 0$
or $\attr$ is already in $\State$.
Otherwise, $\FA$ computes $\sigma \gets \SSign(\mfk, (\attr, F, \vkattr))$,
sends $\sigma$ to $\User$,
and updates
$\State' \coloneqq (\State \| (\attr, F, \transcript \coloneqq (\vkattr, \sigmaenroll, \sigma)))$.
The updated state $\State'$ replaces $\State$ in subsequent executions.
\item
$\User$ proceeds if $\SVerify(\vk, (\attr, F, \vkattr), \sigma) = 1$.
Otherwise, $\User$ aborts.
\item 
$\User$ sets
$\certattr \coloneqq (\attr, F, \vkattr, \sigma)$ and $\secattr \coloneqq \skattr$.
\end{enumerate}

\item[$\SMSSign(\pp, \certattr, \secattr, G, M, w) \to (\tM, \Sigma)$:]
To sign a message $M$,
\
\begin{enumerate}
\item 
Parse $\gpk = (\ek, \vk, \ppnizk)$,
$\certattr = (\attr, F, \vkattr, \sigma)$, $\secattr = \skattr$.
\item 
Compute $(\tM, \op) \coloneqq P(F(M, \attr, w), G(M, \attr, w))$.
\item 
Abort if $\op \notin \OpeningSetSignable$.
Set $\ctop \coloneqq \EEnc(\ek, \op; \rop)$ with $\rop \gets \RandomnessEncrypt$.
\item 
Generate $(\ovk, \osk) \gets \OTSKeyGen(1^\lambda)$
and $\sigmaovk \gets \SSign(\skattr, \ovk)$.
\item 
Execute the prover of $\NIZKAoK$ to generate $\PiGenericSMS$, for $\RelationGenericSMS$
(as in \eqref{equation:relation-g-SMS})
with the statement $(\pp, G, \tM, \ctop, \ovk)$
and witness $(\attr, \vkattr, F, \allowbreak \sigma, M, w, \op, \rop, \sigmaovk)$.
\item 
Compute $\sigmaonetime \gets \OTSSign(\osk, (G, \tM, \ctop, \PiGenericSMS))$.
\item 
Return
$
\bigl(\tM,\; \SigmaGenericSMS\bigr)
= \bigl(\tM,\; (\ctop, \PiGenericSMS, \ovk, \sigmaonetime)\bigr).
$
\end{enumerate}

\end{description}

The relation $\RelationGenericSMS$ certified by $\PiGenericSMS$ is:
\begin{equation}\label{equation:relation-g-SMS}
(\pp, G, \tM, \ctop, \ovk;\;
\attr, \vkattr, F, \sigma, M, w, \op, \rop, \sigmaovk)
\in \RelationGenericSMS
\end{equation}
\[
\iff
\begin{cases}
\SVerify(\vk, (\attr, F, \vkattr), \sigma) = 1,
\\
(\tM, \op) = P(F(M, \attr, w), G(M, \attr, w))
\;\land\; \op \in \OpeningSetSignable,
\\
\ctop = \EEnc(\ek, \op;\, \rop)
\;\land\; \SVerify(\vkattr, \ovk, \sigmaovk) = 1.
\end{cases}
\]

The relation captures three obligations simultaneously:
the signer holds a valid $\FA$-issued certificate over $(\attr, F)$;
the combiner evaluation is correct and the resulting tag is signable;
and $\op$ is honestly encrypted for $\TA$ under $\ek$.

\begin{description}

\item[$\SMSVerify(\pp, G, \tM, \SigmaGenericSMS) \to \{0, 1\}$:]
To verify $\SigmaGenericSMS$:
\begin{enumerate}
\item 
Parse $\gpk = (\ek, \vk, \ppnizk)$
and $\SigmaGenericSMS = (\ctop, \PiGenericSMS, \ovk, \sigmaonetime)$.
\item 
Return $1$ if and only if
$\OTSVerify(\ovk, (G, \tM, \ctop, \PiGenericSMS), \sigmaonetime) = 1$
and $\PiGenericSMS$ is valid for $\RelationGenericSMS$
with public parameter $\ppnizk$ and statement $(\ek, \vk, G, \tM, \ctop, \ovk)$.
\end{enumerate}

\item[$\SMSOpen(\pp, \tsk, G, \tM, \SigmaGenericSMS) \to \OpeningSet$:]
Parse $\SigmaGenericSMS = (\ctop, \PiGenericSMS, \ovk, \sigmaonetime)$
and $\tsk = \dk$.
Return $\op' \coloneqq \EDec(\dk, \ctop)$.

\end{description}
% !TEX root = main.tex
\subsection{Analysis}\label{subsection:generic-SMS-analysis}

\begin{theorem}[Security]\label{theorem:SMS-generic-construction-security}
$\SMSSchemeGeneric$ satisfies correctness, privacy, setup indistinguishability,
and unforgeability
(Definitions~\ref{definition:sms-correctness},
\ref{definition:sms-privacy},
\ref{definition:sms-setup-indistinguishability},
and \ref{definition:sms-unforgeability}),
assuming that $\NIZKAoK$ is a simulation-sound extractable argument of knowledge,
$\SignatureScheme$ is existentially unforgeable,
$\OTSScheme$ is strongly unforgeable,
and $\EncryptionScheme$ is CPA-secure.
\end{theorem}

\begin{proof}[Sketch]
Correctness follows directly from the correctness of the underlying primitives.
We order the properties by logical dependency.

\medskip
\noindent\textbf{Privacy.}
The only signature-dependent information visible to the adversary is
$\ctop$ and $\PiGenericSMS$.
Since the two candidate signers produce the same $\tM$ and the same $\op$
(enforced by the privacy experiment),
$\ctop$ is an encryption of the same value in both cases;
CPA security of $\EncryptionScheme$ ensures $\ctop$ is indistinguishable.
The zero-knowledge property of the $\NIZKAoK$ ensures $\PiGenericSMS$
reveals no information about the witness,
including $\attr$, $F$, and the computation trace.
Together, the two components are jointly indistinguishable across the two signers.

\medskip
\noindent\textbf{Setup indistinguishability.}
$\SMSSimSetup$ is identical to $\SMSSetup$ except that it additionally generates
an extraction trapdoor $\tauext$ (deferred to the full version).
All other components, namely, $\SignatureScheme$, $\OTSScheme$, $\EncryptionScheme$,
are executed identically in the real and simulated setups.
Setup indistinguishability follows from the simulation indistinguishability
of $\NIZKAoK$: any adversary distinguishing the two setups yields
a distinguisher for $\NIZKAoK$'s simulation indistinguishability.

\medskip
\noindent\textbf{Extractability.}\hspace{-0.7pt}
Given a valid forgery $\!(G, \tM, \SigmaGenericSMS\hspace{-0.7pt})$,
$\SMSExtract$ runs $\NIZKAoK$'s extractor on $\PiGenericSMS$
using trapdoor $\tauext$ to recover $(\attr', F', M', w', \op', \rop', \sigmaovk')$.
The simulation-sound extractability~\cite{CRYPTO:ChaLys06}
guarantees that the extracted witness satisfies $\RelationGenericSMS$,
so $\op' \in \OpeningSetSignable$, $\op' = \op$, and $\tildeM = \tM$ hold.

\medskip
\noindent\textbf{Type-1 unforgeability (traceability).}
Suppose $\Adversary$ forges a signature with extracted $(\attr', F') \notin \LCU$.
By extractability, the extracted $\sigma'$ verifies, meaning
$\SVerify(\vk, (\attr', F', \vkattr'), \sigma') = 1$.
Since $(\attr', F')$ was never enrolled by $\FA$ (not in $\LCU$),
no honest $\SMSJoin$ execution produced $\sigma'$,
and $\Adversary$ did not receive $\mfk$ (as $\OMfk$ is restricted).
This yields a forgery against $\SignatureScheme$ under verification key $\vk$.

\medskip
\noindent\textbf{Type-2 unforgeability (non-frameability).}
Suppose $\Adversary$ forges a signature with extracted $(\attr', F') \in \LHU$.
By extractability, $\SVerify(\vkattr', \ovk', \sigmaovk') = 1$,
where $\vkattr'$ is the honest user's verification key recorded in $\State$.
The one-time signature $\sigmaonetime$ on $(G, \tM, \ctop, \PiGenericSMS)$
verifies under $\ovk'$, which was never signed by the honest user's $\skattr'$
(since the forgery is fresh relative to $\LSig$).
This yields a forgery against the existential unforgeability of $\SignatureScheme$
under $\vkattr'$, the honest user's per-user verification key.

The role of $\OTSScheme$ is distinct: strong unforgeability of $\OTSScheme$
prevents an adversary from taking a legitimately obtained signature
$(\ctop, \PiGenericSMS, \ovk, \sigmaonetime)$ and modifying any component
to produce a new valid signature, since $\sigmaonetime$ binds all components
under $\ovk$, which is itself certified under $\vkattr$.

\medskip
Details (\eg, $\SMSSimSetup$ and $\SMSExtract$)
are deferred to \ifcameraready the full version.
\else Appendix~\ref{appendix:SMS-analysis-generic-construction}.
\fi
\qed
\end{proof}
